\documentclass[10pt,a4paper]{article}
\usepackage[T1]{fontenc}\usepackage[utf8]{inputenc}\usepackage[english,italian]{babel}
\usepackage{lmodern,microtype}\usepackage[a4paper,margin=1.85cm,headheight=15pt]{geometry}
\usepackage{enumitem,tabularx,array,longtable,booktabs,xcolor,hyperref,xurl,fancyhdr,titlesec,framed}
\definecolor{ddtadark}{RGB}{28,38,52}\definecolor{ddtablue}{RGB}{42,88,135}\definecolor{ddtagray}{RGB}{244,246,248}\definecolor{ddtagreen}{RGB}{48,111,78}\definecolor{ddtaorange}{RGB}{148,91,25}
\hypersetup{colorlinks=true,linkcolor=ddtablue,urlcolor=ddtablue}\pagestyle{fancy}\fancyhf{}\lhead{DDTA - BA0-R}\rhead{REVIEW-ONLY}\cfoot{\thepage}
\titleformat{\section}{\large\bfseries\color{ddtadark}}{}{0pt}{}\setlist[itemize]{leftmargin=1.45em,itemsep=.18em,topsep=.2em}
\setlength{\parindent}{0pt}\setlength{\parskip}{.32em}\setlength{\emergencystretch}{2em}
\newenvironment{factbox}[1]{\def\FrameCommand{\fcolorbox{ddtablue}{ddtagray}}\MakeFramed{\advance\hsize-2\width\FrameRestore}\noindent\textbf{#1}\par\smallskip}{\endMakeFramed}
\newenvironment{challenge}[1]{\def\FrameCommand{\fcolorbox{ddtaorange}{ddtagray}}\MakeFramed{\advance\hsize-2\width\FrameRestore}\noindent\textbf{#1}\par\smallskip}{\endMakeFramed}

\begin{document}
\begin{center}
{\LARGE\bfseries DDTA - Scheda di lettura compilata}\par
{\large SRC-0039 - BA0-R systems-modeling prior art}\par
{\small REVIEW-ONLY: da approvare prima del caricamento nel repository.}\par
\end{center}
\begin{factbox}{Identita bibliografica e accesso}
\begin{tabularx}{\linewidth}{>{\bfseries}p{3.4cm}X}
Source ID & SRC-0039\\
Titolo & NASA Systems Modeling Handbook for Systems Engineering\\
Autori/ente & National Aeronautics and Space Administration (NASA)\\
Anno & 2025\\
Identificatore & NASA-HDBK-1009A, Revision A, 2025-03-12\\
Accesso & internet\_public\_official\_nasa\_pdf\\
Verification state & verified\_full\_text\_public\\
\end{tabularx}
\end{factbox}
\section{1. Research problem and contribution}
Integrates system modeling with selected NASA systems-engineering processes and work products. The handbook is tool-agnostic and states that, although it uses SysML v1.7 examples, other modeling languages may be applied to its metamodel. It separates modeling language, methodology and framework and uses a metamodel to relate elements across structure, behavior, requirements and parametrics.
\section{2. Starting artifacts and generated representations}
The model captures stakeholder expectations, requirements, structure, behavior, interfaces, item/information flows, measures, verification and validation knowledge. System context captures scope/boundaries, the system of interest, users and external systems. Internal block views show structural interfaces and items flowing across them. Diagrams and tables are generated/extracted from the populated model rather than treated as the canonical knowledge store.
\section{3. Traceability, change, automation and human role}
The handbook treats model planning, conventions, accessibility, roles/responsibilities and model management as explicit planning concerns. It uses traceability relationships across requirements, measures, verification and structure/behavior elements. Rendering choices may omit ports or details for visual simplicity without deleting model content. Tooling supports extraction of views but does not replace modeling methodology or governance.
\section{4. Findings, limitations and relationship with DDTA}
Strong prior art for a methodology-neutral canonical system model with multiple projections. It directly weakens a minimal Actor+Asset+Interaction core because interface, flow, item, structure, behavior and mode are distinct responsibilities. It does not establish that NASA classes should become DDTA BAE metaclasses; its metamodel is tailored to NASA SE work products and SysML practice. It strongly supports Base Analysis as a canonical knowledge representation from which views are derived.
\section{Evidenza utile per BA0/BA1}
\begin{longtable}{p{1.2cm}p{2.0cm}p{4.0cm}p{8.0cm}}
\toprule
Pag. & Ruolo & Sezione & Interpretazione DDTA\\
\midrule
\endhead
5 & foundation & 1.1 Purpose & A modeling metamodel can be conceptually separated from its concrete notation/tool.\\
\addlinespace[.3em]
11 & foundation & 4.3.1.2 Modeling Pillars of SysML & Structure and behavior are distinct system-modeling responsibilities rather than one generic interaction bucket.\\
\addlinespace[.3em]
14 & foundation & 4.3.3 Modeling Framework & Canonical model content should precede view generation.\\
\addlinespace[.3em]
23 & foundation & 8.3 System Context bdd & Boundary/context is a generic systems-modeling responsibility, not inherently a security trust boundary.\\
\addlinespace[.3em]
24 & challenge & 8.4 System Context ibd & Interface and transferred item/flow should not be collapsed into a single generic interaction.\\
\addlinespace[.3em]
45 & foundation & 9 Generating Diagrams and Tables & Views are projections/renderings of canonical knowledge, not independent sources of truth.\\
\addlinespace[.3em]
70 & comparison & Appendix F 3.3 Interfaces & An interface abstraction can be broader than a communication channel.\\
\addlinespace[.3em]
72 & foundation & Appendix F 3.4 Modes of Operations & State/mode is a recurring semantic responsibility that BA1 must test.\\
\addlinespace[.3em]
\bottomrule
\end{longtable}
\section{Bibliography mining / follow-up}
\begin{itemize}
\item OMG SysML v1.7 and v2 specifications
\item Friedenthal, Moore \& Steiner, Practical Guide to SysML
\item NASA OpenSE Cookbook / MBSE Grid references
\end{itemize}
\begin{challenge}{Governance note}
Questa scheda e un ausilio di review. Le future citazioni di tesi devono risolversi alla fonte registrata e a una posizione esatta nell'excerpt ledger approvato. Nessun concetto esterno diventa automaticamente un BAE DDTA.
\end{challenge}
\end{document}
